% =============================================================================
% TIKZ SNIPPETS FOR BEAMER
%
% Ready-to-paste TikZ patterns for academic presentations.
% All snippets use PRIMARY as the accent color — works with the template's
% \definecolor{PRIMARY}{rgb}{...} definition.
%
% Usage: copy the snippet inside a \begin{frame}...\end{frame} block.
%        Wrap in \begin{center}...\end{center} for centered placement.
%        Adjust labels, node distances, and sizes to fit your content.
%
% Requires (already in template): \usepackage{tikz}
% Some snippets also need: \usetikzlibrary{arrows.meta,positioning,fit}
% Add to preamble if missing.
% =============================================================================


% -----------------------------------------------------------------------------
% 1. DAG — Identification diagram
%    Common use: instrument Z → treatment X → outcome Y, with confounder U
% -----------------------------------------------------------------------------
%
% \usetikzlibrary{arrows.meta, positioning}
%
% \begin{tikzpicture}[
%     >=Stealth,
%     node distance = 2.2cm,
%     obs/.style  = {circle, draw, thick, minimum size=1cm, font=\small},
%     unobs/.style= {circle, draw, thick, dashed, minimum size=1cm,
%                    font=\small, text=gray}
%   ]
%   \node[obs]   (Z) {$Z$};
%   \node[obs]   (X) [right=of Z]            {$X$};
%   \node[obs]   (Y) [right=of X]            {$Y$};
%   \node[unobs] (U) [above right=1cm and 1cm of X] {$U$};
%
%   \draw[->]               (Z) -- (X);
%   \draw[->, color=PRIMARY](X) -- (Y)  node[midway, below] {\scriptsize effect};
%   \draw[->, dashed, gray] (U) -- (X);
%   \draw[->, dashed, gray] (U) -- (Y);
%   % Exclusion restriction: Z does not directly affect Y
%   % (add a crossed-out arrow if you want to make this explicit)
% \end{tikzpicture}


% -----------------------------------------------------------------------------
% 2. Event study timeline
%    Common use: visualize pre/post periods around a treatment event
% -----------------------------------------------------------------------------
%
% \usetikzlibrary{arrows.meta}
%
% \begin{tikzpicture}[>=Stealth]
%   % Main axis
%   \draw[->] (-5,0) -- (5,0) node[right] {\small Time};
%
%   % Tick marks and labels
%   \foreach \t/\lbl in {-4/$t-4$, -3/$t-3$, -2/$t-2$, -1/$t-1$,
%                         0/$t=0$,   1/$t+1$,  2/$t+2$,  3/$t+3$} {
%     \draw (\t, 0.12) -- (\t, -0.12);
%     \node[below, font=\scriptsize] at (\t, -0.12) {\lbl};
%   }
%
%   % Pre-period bracket
%   \draw[decorate, decoration={brace, amplitude=6pt}, color=gray]
%     (-4, 0.4) -- (-0.1, 0.4) node[midway, above=6pt, font=\scriptsize, gray]
%     {Pre-period};
%
%   % Post-period bracket
%   \draw[decorate, decoration={brace, amplitude=6pt}, color=PRIMARY]
%     (0.1, 0.4) -- (3, 0.4) node[midway, above=6pt, font=\scriptsize, color=PRIMARY]
%     {Post-period};
%
%   % Treatment line
%   \draw[dashed, thick, color=PRIMARY] (0, -0.5) -- (0, 1.3)
%     node[above, font=\small, color=PRIMARY] {Treatment};
% \end{tikzpicture}
%
% Note: add \usetikzlibrary{decorations.pathreplacing} for the braces.


% -----------------------------------------------------------------------------
% 3. Difference-in-differences 2×2 table
%    Common use: illustrate the DiD logic with four cells
% -----------------------------------------------------------------------------
%
% \begin{tikzpicture}[font=\small]
%   % Column and row labels
%   \node at (1.5, 3.2)  {\textbf{Pre}};
%   \node at (3.5, 3.2)  {\textbf{Post}};
%   \node[rotate=90] at (-0.3, 2.0) {\textbf{Control}};
%   \node[rotate=90] at (-0.3, 0.5) {\textbf{Treated}};
%
%   % Grid
%   \draw (0,0) rectangle (5,3);
%   \draw (2.5,0) -- (2.5,3);   % vertical divider
%   \draw (0,1.5) -- (5,1.5);   % horizontal divider
%
%   % Cell values — replace with actual estimates or labels
%   \node at (1.25, 2.25) {$\bar{Y}_{C,Pre}$};
%   \node at (3.75, 2.25) {$\bar{Y}_{C,Post}$};
%   \node[color=PRIMARY] at (1.25, 0.75) {$\bar{Y}_{T,Pre}$};
%   \node[color=PRIMARY] at (3.75, 0.75) {$\bar{Y}_{T,Post}$};
%
%   % DiD arrow (optional)
%   \draw[->, thick, color=PRIMARY] (2.7, 0.75) -- (4.9, 0.75)
%     node[midway, below, font=\scriptsize] {DiD};
% \end{tikzpicture}


% -----------------------------------------------------------------------------
% 4. Methodology pipeline / flowchart
%    Common use: show a sequence of steps in your empirical or theoretical method
% -----------------------------------------------------------------------------
%
% \usetikzlibrary{arrows.meta, positioning}
%
% \begin{tikzpicture}[
%     >=Stealth,
%     box/.style = {draw, rectangle, rounded corners=3pt,
%                   minimum width=2.4cm, minimum height=0.9cm,
%                   align=center, font=\small, fill=PRIMARY!15, thick},
%     node distance = 0.7cm   % gap between boxes
%   ]
%   \node[box] (s1) {Step 1\\{\scriptsize Description}};
%   \node[box] (s2) [right=of s1] {Step 2\\{\scriptsize Description}};
%   \node[box] (s3) [right=of s2] {Step 3\\{\scriptsize Description}};
%   \node[box] (s4) [right=of s3] {Step 4\\{\scriptsize Description}};
%
%   \draw[->] (s1) -- (s2);
%   \draw[->] (s2) -- (s3);
%   \draw[->] (s3) -- (s4);
%
%   % Optional: highlight one step
%   % \node[box, fill=PRIMARY!50, draw=PRIMARY, thick] (s2) ... ;
% \end{tikzpicture}


% -----------------------------------------------------------------------------
% 5. Regression discontinuity cutoff diagram
%    Common use: illustrate the RD design with a cutoff on a running variable
% -----------------------------------------------------------------------------
%
% \usetikzlibrary{arrows.meta}
%
% \begin{tikzpicture}[>=Stealth, font=\small]
%   % Axes
%   \draw[->] (-4.5, 0) -- (4.5, 0) node[right] {Running variable $x$};
%   \draw[->] (0, -0.3) -- (0,  3.0) node[above] {Outcome $Y$};
%
%   % Cutoff line
%   \draw[dashed, thick, color=PRIMARY] (1, -0.3) -- (1, 2.8)
%     node[above, color=PRIMARY] {Cutoff $c$};
%
%   % Left of cutoff: control group (solid line)
%   \draw[thick, gray]    (-4, 0.8) .. controls (-2, 1.0) .. (0.95, 1.3);
%   % Right of cutoff: treated group (solid line, jump up)
%   \draw[thick, color=PRIMARY] (1.05, 2.0) .. controls (2.5, 2.2) .. (4, 2.4);
%
%   % Jump arrow
%   \draw[<->, thick, color=PRIMARY] (1.3, 1.3) -- (1.3, 2.0)
%     node[midway, right] {\scriptsize $\hat\tau$};
%
%   % Labels
%   \node[gray,  font=\scriptsize] at (-2.5, 0.5) {Control ($x < c$)};
%   \node[color=PRIMARY, font=\scriptsize] at ( 2.8, 1.7) {Treated ($x \geq c$)};
% \end{tikzpicture}
